\section{Introduction}
\label{sec:kappa:introduction}

This is the introduction of the kappa. This work is important~\cite{Alon00Algebraic,Alon99Combinatorial,AR23MCSP}. This problem has been studied in the literature~\cite{AN05Chromatic,AR22PerfectMatchingJournal,BBKMW21SpectralPlanting}.

The thesis is structured as follows. First, the so-called kappa gives a comprehensive summary over the thesis. After that the papers included in this thesis are given. The kappa starts with some preliminaries in \cref{sec:kappa:preliminaries}. \dots